% Chapter 2 -- Ohm's law, power and the simulator
\chapter{Ohm's law, power and the simulator}
\label{ch:ohm}

Chapter~1 gave you the nouns: voltage, current, resistance, schematic, node.
This chapter gives you the two verbs — compute and simulate — and settles the
debt from Example~1.2 by actually choosing the LED's resistor. When you finish
it you should be able to: use $V=IR$ and $P=VI$ in every algebraic form,
check whether a real resistor survives the power you ask of it, read the
headline numbers on a component datasheet, write and run your first
\emph{ngspice} netlist (the simulator every analogue chapter uses), and say
honestly when Ohm's law applies and when it does not.

\section{Ohm's law: a definition with a job}
\label{sec:ohmslaw}

In Chapter~1 resistance was introduced as $R=V/I$ — volts per amp, measured in
ohms\index{ohm}. Turned around, that is \emph{Ohm's law}\index{Ohm's law}:
\[
  V = I\,R \qquad\text{and the two shadows}\qquad I = \frac{V}{R},
  \quad R = \frac{V}{I}.
\]
For the components this law describes — resistors, wire, most fixed loads —
it says the current is \emph{proportional} to the voltage. Double the push,
double the flow; the ratio never changes. That constancy is the content: a
resistor is a part whose $V$-$I$ graph is a straight line through the origin,
and $R$ is its slope\index{V-I characteristic}. Components for which that is
true are called \emph{linear}\index{linear element} or \emph{ohmic}. Most of
the interesting ones are not — Section~\ref{sec:nonlinear} is about that.

The law earns its keep on one more fact we will prove properly in
Chapter~3: in a single loop, the source's voltage is \emph{spent} across the
elements in the path — the shares always add back to the supply. If the
\qty{5}{\volt} battery in Figure~1.3 pushes and the LED keeps about
\qty{2}{\volt} for itself, the resistor is left holding the other
\qty{3}{\volt}. That sentence is the entire solution to the problem we posed
in Chapter~1 and could not yet solve.

\begin{example}[The LED question, answered twice]
\label{ex:ledsizing}
\emph{Problem.} Drive an LED from \qty{5}{\volt} and, separately, from
\qty{3.3}{\volt}. Datasheet headline: $V_F \approx 2$\,V at the intended
current, maximum $I_F = 30$\,mA. Target \qty{15}{\milli\ampere}.\par
\emph{Strategy.} The resistor does two jobs: it drops whatever voltage the
LED does not take, and Ohm's law then converts that drop into the current
everything in the loop shares. Compute $R$, pick the nearest standard E24
value, then check the corners: source $\pm 5\%$, resistor $\pm 5\%$,
$V_F$ anywhere in $1.8$--$2.2$\,V.\par
\emph{Solution, \qty{5}{\volt}.} $R = (5-2)/0.015 = \qty{200}{\ohm}$ — itself
an E24 value, so the nominal build gives exactly \qty{15}{\milli\ampere}.
Corners:
\[
  I_{\min} = \frac{5\times0.95 - 2.2}{200\times1.05}
           = \frac{2.55}{210} = 12.1\ \mathrm{mA},
  \qquad
  I_{\max} = \frac{5\times1.05 - 1.8}{200\times0.95}
           = \frac{3.45}{190} = 18.2\ \mathrm{mA}.
\]
Every combination is legal: nowhere near 30\,mA, brightness variation
invisible.\par
\emph{Solution, \qty{3.3}{\volt}.} $R = (3.3-2)/0.015 = \qty{86.7}{\ohm}
\rightarrow$ choose E24 \qty{91}{\ohm} (nominal 14.3\,mA). Corners:
\[
  I_{\min} = \frac{3.3\times0.95-2.2}{91\times1.05} = 9.8\ \mathrm{mA},
  \qquad
  I_{\max} = \frac{3.3\times1.05-1.8}{91\times0.95} = 19.3\ \mathrm{mA}.
\]
\emph{Notice:} the lower the headroom over $V_F$, the wider the proportional
spread — \qty{3.3}{\volt} designs run 10--19\,mA where \qty{5}{\volt} ran
12--18. Sensitivity analysis is not pedantry; it is how you find the
designs that fail only on Tuesdays.
\end{example}

Example~\ref{ex:ledsizing} ends with a part number: ``\qty{200}{\ohm}.''
The schematic knows that; the beige cylinder in your hand does not carry a
number at all — it carries paint. Figure~\ref{fig:colorcode} shows how to
read it.

\begin{figure}[htbp]
\centering
% Figure 2.1: reading a 4-band resistor
\begin{tikzpicture}[scale=1.0,font=\small]
  % resistor body
  \draw[thick] (-4.2,0) -- (-3.1,0);
  \draw[thick] (3.1,0) -- (4.2,0);
  \draw[thick,fill=graybg,rounded corners=2pt] (-3.1,-0.55) rectangle (3.1,0.55);
  % bands: red red brown gold = 220 ohm
  \fill[red!75!black]     (-2.55,-0.55) rectangle (-2.05,0.55);
  \fill[red!75!black]     (-1.35,-0.55) rectangle (-0.85,0.55);
  \fill[brown!70!black]   ( 0.05,-0.55) rectangle ( 0.55,0.55);
  \fill[yellow!60!orange!70] ( 2.15,-0.55) rectangle ( 2.65,0.55);
  % callouts: stagger two depths to avoid collisions
  \draw[->,slate] (-2.3,-0.6) -- (-2.3,-1.1);
  \node[below,align=center,font=\scriptsize] at (-2.3,-1.1) {1st digit\\[-2pt]red = 2};
  \draw[->,slate] (-1.1,-0.6) -- (-1.1,-1.75);
  \node[below,align=center,font=\scriptsize] at (-1.1,-1.75) {2nd digit\\[-2pt]red = 2};
  \draw[->,slate] (0.3,-0.6) -- (0.3,-1.1);
  \node[below,align=center,font=\scriptsize] at (0.3,-1.1) {multiplier\\[-2pt]brown = $\times 10^{1}$};
  \draw[->,slate] (2.4,-0.6) -- (2.4,-1.75);
  \node[below,align=center,font=\scriptsize] at (2.4,-1.75) {tolerance\\[-2pt]gold = $\pm 5\%$};
  \node[align=center,font=\small] at (0,-3.1)
    {$2\,2\,\times 10^{1} = \qty{220}{\ohm}\pm5\%$};
  % read direction note
  \node[slate,align=left,font=\scriptsize,anchor=west] at (4.5,-1.1)
    {read digits toward the\\[-2pt]multiplier; the tolerance\\[-2pt]band sits alone at the end};
\end{tikzpicture}

\caption{The 4-band resistor code: two significant digits, a power-of-ten
multiplier, and a tolerance band set apart at the end. Read digits toward
the multiplier. Notice that the bands encode the \emph{nominal} value only —
the true part lives anywhere inside the tolerance band, which is why
Example~\ref{ex:ledsizing} computed corners, not just the centre.}
\label{fig:colorcode}
\end{figure}

\begin{exercise}[easy]
Compute: (a) $I$ for $\qty{9}{\volt}$ across $\qty{330}{\ohm}$;
(b) $R$ needed for \qty{10}{\milli\ampere} from \qty{3.3}{\volt} (no LED);
(c) the voltage a \qty{1}{\kilo\ohm} resistor drops at \qty{2}{\milli\ampere}.
\end{exercise}

\section{Power: $P=VI$ and the quarter-watt ceiling}
\label{sec:power}

Voltage is energy per charge; current is charge per second. Multiply them
and the units do the talking:
\[
  P = V\,I, \qquad \text{1 watt} = \text{1 joule per second}.
  \index{power|textbf}\index{watt}
\]
Every element in a circuit either supplies this power (sources, negative in
the bookkeeping convention) or spends it (everything else). A resistor spends
it as heat — substituting $V=IR$ gives the two forms you will use constantly:
\[
  P = I^2 R = \frac{V^2}{R}. \index{power!dissipation}
\]
The second form has a counterintuitive lesson: \emph{less} resistance across a
fixed voltage means \emph{more} dissipation. A $\qty{1}{\ohm}$ resistor across
\qty{5}{\volt} dissipates 25\,W — a space heater, if it survives.

Components are sold with a \emph{power rating}\index{power rating}: the
continuous dissipation their package can shed before it overheats. The
jellybean through-hole resistor of this book is rated
\qty{0.25}{\watt} (a quarter watt). Ratings are not suggestions —
Example~\ref{ex:rating} shows the failure mode, and Section~\ref{sec:safety2}
the safety one.

\begin{example}[Can a 1/4 W resistor carry 100 mA?]
\label{ex:rating}
\emph{Problem.} A design wants \qty{100}{\milli\ampere} through a
\qty{47}{\ohm} resistor. Is the standard \qty{0.25}{\watt} part
adequate?\par
\emph{Strategy.} The rating is about heat, so compute the dissipation with the
form that uses the given quantities: $I$ and $R$ are known, so $P=I^2R$.\par
\emph{Solution.} $P = (0.1)^2 \times 47 = 0.47$\,W $> 0.25$\,W. \textbf{Fails}
— by nearly a factor of two. A \qty{0.5}{\watt} part barely passes
($0.47 < 0.5$), and engineering practice is to run resistors at no more than
half to two-thirds of their rating (the \emph{derating}\index{derating}
habit: ratings assume ideal mounting, cool air, and a young part). So the
correct order is a \qty{1}{\watt} resistor — or better, a design that does not
need 100\,mA.\par
\emph{Notice:} nobody warned you until you computed. Power ratings are
invisible on the schematic; checking them is a habit, not a step.
\end{example}

\begin{exercise}[medium]
Repeat the check for Example~\ref{ex:ledsizing}: the \qty{200}{\ohm} resistor
at \qty{15}{\milli\ampere} dissipates how much? Does \qty{0.25}{\watt} suffice
with derating margin? Then the worst-corner case: \qty{190}{\ohm} at
18.2\,mA.
\end{exercise}

\section{When the line bends: nonlinear loads}
\label{sec:nonlinear}

Here is the sentence to un-learn early: ``the LED's resistance is
$2\,\mathrm{V}/20\,\mathrm{mA} = \qty{100}{\ohm}$.'' That number exists, but
it is not a resistance — it is one point on a curve, and the curve is nothing
like a line.

\begin{figure}[htbp]
\centering
% Figure 2.1: resistor V-I line vs LED I-V curve (diode data = real ngspice sweep,
% figures/ch02/ch02_diode_iv.dat, produced by code/ch02_diode_iv.cir)
\begin{tikzpicture}
\begin{axis}[
  width=12cm,height=6.8cm,
  xlabel={voltage across the part (V)},
  ylabel={current through it (mA)},
  xmin=0,xmax=3,ymin=0,ymax=20,
  grid=both,grid style={rulegray!40},
  legend style={at={(0.03,0.97)},anchor=north west,draw=none,font=\small,fill=none},
  tick label style={font=\small},label style={font=\small},
]
% resistor: I = V/200 ohm -> mA = 1000*V/200 = 5*V
\addplot[domain=0:3,samples=2,thick,ink] {5*x};
\addlegendentry{\qty{200}{\ohm} resistor: $I=V/R$, a straight line through 0}
% LED curve from real ngspice data (amps -> mA)
\addplot[thick,accent] table[x index=0,y index=1] {figures/ch02/ch02_diode_iv_ma.dat};
\addlegendentry{LED (diode model): nothing until $\sim$2 V, then a cliff}
% guide lines
\draw[dashed,slate] (axis cs:2,0) -- (axis cs:2,20);
\node[slate,font=\scriptsize,anchor=west] at (axis cs:2.03,18.5) {$V_F \approx 2$ V};
\end{axis}
\end{tikzpicture}

\caption{A resistor's $V$-$I$ characteristic against an LED's, from a real
ngspice sweep of the diode model used in this chapter's netlists. Notice the
two shapes this chapter keeps separate: the resistor's straight line through
zero (Ohm's law is its \emph{definition}) and the LED's dead-flat-then-cliff
curve (Ohm's law has no foothold at all). Below $V_F$ the LED conducts almost
nothing; above it, the current explodes — which is why the series resistor is
not optional.}
\label{fig:vi}
\end{figure}

A \emph{nonlinear} element\index{nonlinear element} is one whose $V$-$I$
relationship is not a line through the origin. The LED is the tamest case:
a \emph{forward voltage}\index{forward voltage} $V_F$ it holds nearly
constant (about \qty{2}{\volt} for the red indicator type this book uses),
and a current that rises exponentially once you exceed it. $V=IR$ cannot
describe that — the ratio $V/I$ is different at every current. What survives
is the \emph{method}: the LED's curve still passes through every point the
circuit forces it through; ``about \qty{2}{\volt}'' is simply the datasheet's
compression of that curve into one usable number. Chapter~5 unpacks the
device physics; for now, treat $V_F$ as a fact with a tolerance.

\begin{remark}
This is also why the LED \emph{needs} its resistor. A perfect \qty{2}{\volt}
source across a 2-V LED is a knife-edge: a degree of temperature or a
millivolt of mismatch and the exponential current is unbounded. The resistor
converts the knife-edge into a gentle slope (Figure~\ref{fig:sweep}): above
the knee the current is set by $(V_{\mathrm{supply}}-V_F)/R$ — Ohm's law
reappearing once the diode's share is paid.
\end{remark}

\section{Resistor networks, concept only}
\label{sec:networks}

Resistors almost never appear alone. Two facts carry you until Chapter~3's
formal treatment:
\begin{itemize}
\item \textbf{In series}\index{series}, the same current threads all of them,
and their resistances add: $R_{\mathrm{eq}} = R_1 + R_2 + \cdots$ This is why
the divider in Figure~\ref{fig:predict} splits the supply: two equal shares
of the total resistance drop equal shares of the voltage.
\item \textbf{In parallel}\index{parallel}, they share the voltage and split
the current — two \qty{1}{\kilo\ohm} resistors across the same nodes behave
as \qty{500}{\ohm}. More paths, less total opposition.
\end{itemize}
Chapter~3 turns these into Kirchhoff's laws and the divider formula; today
you need only the intuition and, in the exercises, a netlist that checks it.

\section{Reading a datasheet's headline numbers}
\label{sec:datasheet}

Datasheets are long; the skill is knowing which three lines you actually
need for a given decision. For an LED, these:

\begin{table}[htbp]
\centering\small
\begin{tabularx}{\linewidth}{lYl}
\toprule
\thead{Line} & \thead{What it means} & \thead{Typical (5 mm red LED)}\\
\midrule
$V_F$ (forward voltage) & voltage it keeps for itself at working current & 1.8--2.2\,V, ``typ 2.0''\\
$I_F$ max & current you must stay under, always & 20--30\,mA\\
$P_D$ (power dissipation) & heat ceiling, $V_F \times I_F$ & $\sim$100\,mW\\
\midrule
viewing angle, $\lambda$ & optics, not electricity & 20--30$^\circ$, 620--630\,nm\\
\bottomrule
\end{tabularx}
\caption{The three electrical numbers that decide an LED circuit, and the
optics rows you may ignore for now. Notice ``typ'' vs ``max'': \emph{typical}
is marketing's average; \emph{max/min} are the corners your design must
survive — Example~\ref{ex:ledsizing} used the $V_F$ range, not the typ
value, and that choice is why its corners computed to something honest.}
\label{tab:datasheet}
\end{table}

The discipline generalises: for every component in this book, the first
datasheet pass asks only \emph{what does it want} (ratings to satisfy) and
\emph{what does it do to you} (ratings to stay under). Everything else is a
second read.

\section{Your first ngspice netlist}
\label{sec:ngspice}

All the arithmetic so far assumed a \qty{2}{\volt} LED. How wrong is that
assumption — 1\%, 30\%? The way engineers ask such questions without burning
parts is \emph{simulation}\index{simulation}: a program that builds the
equations of your circuit and solves them exactly as you would, minus the
arithmetic slips. This book uses \emph{ngspice}\index{ngspice}, the free,
open-source descendant of Berkeley SPICE — the same engine underneath the
expensive packages.

First, proof it is installed — the Chapter~1 script again, now with
\texttt{ngspice} present (Listing~\ref{lst:checktools}).

\begin{lstlisting}[language={},numbers=none,caption={Real output of
\texttt{sh code/check\_tools.sh} on this machine after \texttt{brew install
ngspice}: the line that matters reads \emph{found}. (The script's version
grab was also fixed: ngspice prints a decorative banner first, so it now
reports the first line containing a digit.)},label={lst:checktools02}]
MISSING  iverilog               Verilog simulation (ch11+)
MISSING  verilator              lint / fast simulation (ch12+)
MISSING  yosys                  synthesis (ch13+)
MISSING  nextpnr-himbaechel     place-and-route, Gowin (ch13+)
MISSING  gowin_pack             bitstream packer, Gowin (ch13+)
MISSING  openFPGALoader         programming the board (ch13+)
MISSING  gtkwave                viewing waveforms (ch11+)
found    ngspice                ** ngspice-47 : Circuit level simulation program
MISSING  kicad-cli              PCB work (ch18)
found    python3                Python 3.9.6
\end{lstlisting}

A SPICE \emph{netlist}\index{netlist} is a text file that is to a schematic
what a program is to pseudocode. Listing~\ref{lst:ledcir} is the entire
simulation for the circuit of Figure~1.3 — and Listing~\ref{lst:ledop} is
its real output on this machine.

\begin{lstlisting}[language={},caption={\texttt{code/ch02\_led.cir} -- every
line earns its place: a comment, one source, one resistor, one diode, a
model, an analysis request, \texttt{.end}.},label={lst:ledcir}]
* ch02_led.cir -- operating point of the 5 V, R, LED circuit
* Worked example 2.1 target: about 15 mA through the LED
V1 1 0 DC 5
R1 1 2 200
D1 2 0 LEDMOD
.model LEDMOD D (IS=1e-19 N=2 BV=100)
.op
.end
\end{lstlisting}

Read it like ngspice does. The first line is always a title/comment. Each
element line is \emph{name, node+, node-, value}: \texttt{V1 1 0 DC 5} puts a
\qty{5}{\volt} source between nodes 1 and 0; \texttt{R1 1 2 200} puts
\qty{200}{\ohm} between nodes 1 and 2; \texttt{D1 2 0 LEDMOD} puts a diode
from node 2 to ground, using the parameter set named \texttt{LEDMOD} that the
\texttt{.model} line defines (its \texttt{IS} and \texttt{N} set the curve you
saw in Figure~\ref{fig:vi}). \emph{Node 0 is ground, always} — the reference
every voltage is quoted against. Finally \texttt{.op} asks for the DC
\emph{operating point}\index{operating point}: solve the circuit once, no
sweeps.

\begin{lstlisting}[language={},numbers=none,caption={Real output of
\texttt{ngspice -b code/ch02\_led.cir} (model parameter tables omitted; the
numbers shown are verbatim). Three things to find: the LED's actual $V_F$,
the loop current, and the sign of the source current.},label={lst:ledop}]
        Node                                  Voltage
        ----    -------
        V(2)                             2.045101e+00
        V(1)                             5.000000e+00

        Source  Current
        ------  -------

        v1#branch                       -1.47745e-02

 Diode: Junction Diode model
         vd                2.0451
         id             0.0147745

 Resistor: Simple linear resistor
         i             0.0147745
         p             0.0436571

 Vsource: Independent voltage source
         i            -0.0147745
         p            -0.0738725
\end{lstlisting}

Three readings, all loaded. \textbf{(1)} $v(2)=2.045$\,V — the model's $V_F$,
not the assumed 2.0 — and $i = 14.77$\,mA, not 15. The 1.5\% gap is the
difference between an assumption and a model. \textbf{(2)} Every element
reports the same current — the single loop really does share one flow.
\textbf{(3)} The source's current is \emph{negative}: SPICE measures current
into the positive terminal, and a source delivers current, so its number
comes out negative — convention, not error. Its power is likewise negative:
it is supplying the $+73.9$\,mW the resistor ($+43.7$\,mW) and LED (the
rest) spend.

\begin{figure}[htbp]
\centering
% Figure 2.3: what a netlist is, and where the output comes from
\begin{tikzpicture}[font=\small,>=stealth,
  box/.style={draw,rounded corners=1pt,inner sep=6pt,align=left},
  lbl/.style={font=\scriptsize\itshape,text=slate}]
  % left: the netlist
  \node[box,fill=graybg,text width=4.6cm] (nl) {
    \ttfamily\scriptsize
    * ch02\_led.cir\\[1pt]
    V1 1 0 DC 5\\
    R1 1 2 200\\
    D1 2 0 LEDMOD\\
    .model LEDMOD D (...)\\
    .op\\
    .end};
  % middle: the simulator
  \node[box,fill=accentsoft,anchor=west,align=center] (sim) at ($(nl.east)+(0.9cm,0)$)
    {\textbf{ngspice}\\[-2pt]\scriptsize builds equations\\[-2pt]\scriptsize solves them};
  % right: the output
  \node[box,fill=graybg,text width=4.4cm,anchor=west] (out) at ($(sim.east)+(0.9cm,0)$) {
    \ttfamily\scriptsize
    v(2) = 2.0451e+00\\
    v1\#branch = -1.47745e-02\\[1pt]
    \rmfamily\scriptsize (volts and amps,\\[-2pt]\scriptsize\ttfamily every node and element)};
  \draw[->,thick] (nl) -- (sim) node[midway,above,lbl]{text in};
  \draw[->,thick] (sim) -- (out) node[midway,below,lbl]{numbers out};
  % annotations under netlist
  \node[lbl,anchor=north west,text width=4.6cm] at ($(nl.south west)+(0,-4pt)$)
    {a text file: one line per element, first letter = type (V, R, D), numbers = node names, node 0 = ground};
\end{tikzpicture}

\caption{The simulation loop in one picture: you write a text file of nodes
and elements, ngspice builds and solves the circuit equations, and it reports
every node's voltage and every element's current. The schematic and the
netlist are the same circuit in two languages — matching them is a skill the
exercises drill.}
\label{fig:flow}
\end{figure}

\begin{figure}[htbp]
\centering
% Figure 2.5: the predict-the-output checkpoint -- two-resistor divider with blanks
\begin{circuitikz}[scale=1.0,transform shape]
  \draw (0,0) node[ground]{};
  \draw (0,0) to[battery1,l_=\qty{9}{\volt},invert] (0,7.2);
  \draw (0,7.2) -- (3.4,7.2) node[midway,above]{\scriptsize node 1};
  \draw (3.4,7.2) to[R,l_=\qty{330}{\ohm}] (3.4,4.8);
  \draw (3.4,4.8) -- (4.6,4.8) node[right]{\scriptsize node 2};
  \draw (3.4,4.8) to[R,l_=\qty{330}{\ohm}] (3.4,0) node[ground]{};
  \draw (3.4,0) -- (0,0);
  % blanks
  \node[draw,fill=graybg,rounded corners=1pt,align=center,font=\small] at (7.0,6.0)
    {$v(2)=\ ?\ \mathrm{V}$};
  \node[draw,fill=graybg,rounded corners=1pt,align=center,font=\small] at (7.0,3.6)
    {$i=\ ?\ \mathrm{mA}$};
  \node[slate,font=\scriptsize,align=center] at (3.0,-1.1)
    {predict first — then run the three-line netlist and compare};
\end{circuitikz}

\caption{Checkpoint~\ref{cp:predict}: cover the answer, compute both blanks
by hand, then run the netlist. The point of the exercise is not the numbers —
it is proving to yourself that two equal resistors split a supply evenly.}
\label{fig:predict}
\end{figure}

\begin{checkpoint}[Predict before you run]
\label{cp:predict}
For Figure~\ref{fig:predict} — \qty{9}{\volt}, two \qty{330}{\ohm} resistors
in series — write down $v(2)$ and the loop current before simulating. Then
write the netlist (three element lines) and run \texttt{.op}. Expect:
$v(2)=\qty{4.5}{\volt}$, $i = 9/660 \approx \qty{13.6}{\milli\ampere}$.
\end{checkpoint}

\begin{exercise}[medium]
Write the netlist for Figure~\ref{fig:predict} yourself (title, three
elements, \texttt{.op}, \texttt{.end}), run it, and paste the three numbers
you predicted. If they disagree, find whether the error is in your arithmetic
or your netlist — both happen.
\end{exercise}

\section{Why we simulate before we build}
\label{sec:whysim}

The honest answer is not ``it is easier'' — it is that a simulator makes
\emph{asking questions} free. Changing $R$ from 200 to 100 costs an edit and
a keystroke; on a breadboard it costs a resistor and a resteadied hand. The
sweep commands below are the cheapest form of one question every designer
asks: \emph{what does this circuit do across a range of conditions?}

\begin{lstlisting}[language={},caption={\texttt{code/ch02\_sweep.cir} -- a
\texttt{.dc} sweep walks the source from 0 to 5 V and records the loop
current at each point.},label={lst:sweepcir}]
* ch02_sweep.cir -- dc sweep: LED-circuit current vs supply voltage
V1 1 0 DC 0
R1 1 2 200
D1 2 0 LEDMOD
.model LEDMOD D (IS=1e-19 N=2 BV=100)
.control
dc v1 0 5 0.05
let iload = -i(v1)
set filetype=ascii
wrdata ch02_sweep.dat iload
.endc
.end
\end{lstlisting}

The \texttt{.control} block runs the \texttt{.dc} analysis: sweep \texttt{v1}
from 0 to 5 in steps of 0.05, create a vector for the current (negated so it
reads positive — the source-current sign convention of
Listing~\ref{lst:ledop}), and write the 101 points to a data file.
Listing~\ref{lst:sweepdat} shows real lines from it, and
Figure~\ref{fig:sweep} plots the whole run.

\begin{lstlisting}[language={},numbers=none,caption={Real excerpt of
\texttt{figures/ch02/ch02\_sweep.dat} (voltage, then current in amps): first
line, the knee region, and the endpoint. The endpoint's 14.77 mA is
\emph{exactly} the \texttt{.op} answer — same equations, same result.},
label={lst:sweepdat}]
 0.00000000e+00 -1.60434623e-36
 ...
 1.95000000e+00  4.35823416e-04
 2.00000000e+00  6.02192794e-04
 2.05000000e+00  7.83622176e-04
 ...
 5.00000000e+00  1.47744954e-02
\end{lstlisting}

\begin{figure}[htbp]
\centering
% Figure 2.4: the real .dc sweep (data = ngspice output, figures/ch02/ch02_sweep.dat)
\begin{tikzpicture}
\begin{axis}[
  width=12cm,height=6.8cm,
  xlabel={supply voltage $V_1$ (V)},
  ylabel={circuit current (mA)},
  xmin=0,xmax=5,ymin=0,ymax=17,
  grid=both,grid style={rulegray!40},
  tick label style={font=\small},label style={font=\small},
]
\addplot[thick,accent] table[x index=0,y index=1] {figures/ch02/ch02_sweep_ma.dat};
% annotations
\draw[dashed,slate] (axis cs:2.045,0) -- (axis cs:2.045,17);
\node[slate,font=\scriptsize,anchor=south west] at (axis cs:2.1,1) {LED turns on, $V_F\approx2.05$ V};
\filldraw[ink] (axis cs:5,14.7745) circle (2pt);
\node[font=\scriptsize,anchor=east,align=right] at (axis cs:4.9,15.6)
  {the \texttt{.op} point you ran:\\[-2pt]$(5\,\mathrm{V},\,14.77\,\mathrm{mA})$};
\node[font=\scriptsize,anchor=west,align=left,slate] at (axis cs:0.15,12.5)
  {above the knee the line is straight:\\[-2pt]slope $=1/R$ — Ohm's law\\[-2pt]re-emerging once $V_F$ is paid};
\end{axis}
\end{tikzpicture}

\caption{The full \texttt{.dc} sweep, plotted from the real simulation output.
Notice the two regimes: dead flat below $\sim$2 V (the LED is off and almost
nothing flows anywhere), then a straight ramp — the resistor has taken
control. The marked point is the \texttt{.op} you ran: one point of this
curve computed a different way, agreeing to the last digit.}
\label{fig:sweep}
\end{figure}

\begin{example}[One circuit, three ways]
\label{ex:threeways}
\emph{Problem.} What current does the Figure~1.3 circuit actually draw at
\qty{5}{\volt}?\par
\emph{Strategy.} Compute it three independent ways — assumption, model,
curve — and require them to agree within stated tolerances. When estimates
from different machinery agree, the answer is robust; when they do not, the
disagreement tells you which assumption is weakest.\par
\emph{Solution.}
\begin{enumerate}
\item \textbf{By hand} (the $V_F=2$ assumption): $I = (5-2)/200 =
\qty{15}{\milli\ampere}$.
\item \textbf{By \texttt{.op}} (a model of the real diode curve):
\qty{14.77}{\milli\ampere} — the model's $V_F$ came out 2.045 V, and
$(5-2.045)/200 = 14.77$\,mA exactly, because SPICE just does the same
arithmetic with a sharper input.
\item \textbf{By \texttt{.dc} sweep} (the whole behaviour): Figure~\ref{fig:sweep}
returns the same \qty{14.77}{\milli\ampere} at $V=5$ and shows the number is
on the line's flat, boring stretch — no cliff-edge nearby.
\end{enumerate}
\emph{Notice:} the hand answer was 1.5\% high — inside any tolerance the
circuit will ever see. Simulation did not replace the estimate; it
\emph{priced} it. That is what the simulator is for in this book.
\end{example}

\section{Safety}
\label{sec:safety2}

Chapter~1's rules stand: low voltage only, USB power for everything here.
This chapter adds one line.

\begin{rulebox}[Resistors get hot — check the rating]
A resistor dissipating over its rating does not fail politely; it browns,
smells, and can burn a finger or a board. Before powering anything you
build, run the $P=I^2R$ or $V^2/R$ check of Example~\ref{ex:rating} on the
stressed parts. In this chapter all dissipation is milliwatts — the rule is
training for when it is not.
\end{rulebox}

And a quieter one: \textbf{USB power only} for these experiments. A USB port
is current-limited by design (typically a few amps, and it will simply shut
off or sag rather than push dangerous current into a short). Do not
substitute a bare bench supply or a larger battery until you have a reason —
the limiting is the safety.

\section{Common mistakes}
\label{sec:mistakes2}

\begin{description}
\item[Applying $V=IR$ to nonlinear loads.] The LED has no resistance in the
Ohm's-law sense: its $V/I$ ratio changes at every point (Figure~\ref{fig:vi}).
The fix is not a bigger formula, it is the right model — a $V_F$ fact plus a
series resistor that re-linearises the loop.
\item[Ignoring power ratings until the resistor browns.] Schematics do not
draw heat. A \qty{47}{\ohm} resistor at 100\,mA is a 0.47\,W part inside a
0.25\,W package, and nothing on the schematic page will warn you.
Example~\ref{ex:rating} is the three-line check; run it every time.
\item[Treating a \texttt{.op} number as exact.] Simulation solves the model
you gave it — not the part in your hand. The model said 14.77\,mA; a real
board with a $\pm$5\% resistor and a loose $V_F$ range delivers anywhere in
Example~\ref{ex:ledsizing}'s computed band. Simulate for behaviour and
estimates; measure for truth.
\item[Netlist nodes that don't match the schematic.] Node numbers in a
netlist are names, not positions — but they must connect the same things the
wires do. Every missing wire in a netlist is a component floating in space;
every shared number is an invisible solder joint. Diff the node list against
the schematic before running.
\end{description}

\section{Chapter summary}
\label{sec:ch2summary}

\begin{table}[htbp]
\centering\small
\begin{tabularx}{\linewidth}{llY}
\toprule
\thead{Quantity / idea} & \thead{Form} & \thead{The one-line version}\\
\midrule
Ohm's law & $V=IR$ & defines $R$; holds for linear parts only\\
Power & $P=VI=I^2R=V^2/R$ & energy per second; resistors spend it as heat\\
Power rating & e.g.\ \qty{0.25}{\watt} & package's heat ceiling; derate 2$\times$\\
Nonlinear (LED) & $V_F$, $I_F$ max & no $R$; a cliff curve + a fact from the datasheet\\
Series / parallel & add / share & series: one current, voltages split; parallel: one voltage, currents split\\
Datasheet & typ vs max & satisfy its wants, respect its limits\\
Netlist & name node+ node- value & text schematic; node 0 is ground\\
\texttt{.op} & operating point & one solved state: every $v$ and every $i$\\
\texttt{.dc} & sweep & solve repeatedly over a range — the free question\\
Sign of source $i$ & negative & SPICE convention: current into the $+$ terminal\\
\midrule
Safety & — & rating check first; USB's current limit is a feature\\
\bottomrule
\end{tabularx}
\caption{Chapter 2 on a page: two formulas, one safety habit, and a
simulator that answers what-if questions for free.}
\label{tab:ch2summary}
\end{table}

\section*{Exercises}
\addcontentsline{toc}{section}{Exercises}
\label{sec:ch2exercises}

\begin{exercise}[easy]
(a) Current through \qty{470}{\ohm} at \qty{5}{\volt}. (b) Voltage across
\qty{2.2}{\kilo\ohm} at \qty{3}{\milli\ampere}. (c) $R$ for \qty{8}{\milli\ampere}
from \qty{9}{\volt}.
\end{exercise}

\begin{exercise}[easy]
Dissipation of: \qty{330}{\ohm} at 13.6\,mA; \qty{100}{\ohm} at 5\,V;
\qty{1}{\kilo\ohm} at 20\,mA. Which exceed \qty{0.25}{\watt}?
\end{exercise}

\begin{exercise}[easy]
A datasheet lists $V_F$ typ 2.0\,V, max 2.4\,V at $I_F$=20\,mA, and
$I_F$(abs max)=30\,mA. Which number do you design to, and which do you fear?
\end{exercise}

\begin{exercise}[easy]
In the netlist \texttt{R1 3 7 470; D2 7 0 LEDMOD}, name each element, the
nodes it joins, and where ground is.
\end{exercise}

\begin{exercise}[medium — derive]
\label{ex:derive-power}
Starting from $P=VI$ and $V=IR$, derive $P=I^2R$ and $P=V^2/R$. Then derive
the useful inverse: for a power ceiling $P$ at fixed $V$, the \emph{smallest}
allowed resistance is $R \ge V^2/P$ — watch which way the inequality flips
when you clear a fraction.
\end{exercise}

\begin{exercise}[medium]
Redo Example~\ref{ex:ledsizing} at \qty{5}{\volt} but with $R=\qty{220}{\ohm}$
(the more common bin). Compute the nominal current and the tolerance band;
decide whether it is still a good choice.
\end{exercise}

\begin{exercise}[medium]
Write (don't copy) a netlist for a \qty{9}{\volt} source and two
\qty{1}{\kilo\ohm} resistors in series; predict $v(2)$ and $i$; run
\texttt{.op} and compare.
\end{exercise}

\begin{exercise}[medium]
The kit resistor red-red-brown-gold is nominally \qty{220}{\ohm}$\pm$5\%.
Give the true-value range. A second resistor marked yellow-violet-red-gold:
nominal and range.
\end{exercise}

\begin{exercise}[hard]
Two LEDs in parallel share one series resistor on 5\,V. Explain why this is
a trap: the parts are never identical, so one hogs current, gets hotter,
takes even more current — a runaway split. What does a netlist with two
slightly different \texttt{IS} values show? (Run it.)
\end{exercise}

\begin{exercise}[hard — design]
The supply is unknown, somewhere between 3.3 and 5 V. Choose a resistor and
justify it so the LED current stays within 10--20\,mA across the whole supply
range with $V_F \in [1.8, 2.2]$, or prove the spec is impossible.
\end{exercise}

\subsection*{Hints and selected answers}
\addcontentsline{toc}{subsection}{Hints and selected answers}

\begin{description}
\item[Easy 2.1] 10.6\,mA; 6.6\,V; \qty{1.125}{\kilo\ohm} (E24: \qty{1.1}{\kilo\ohm}).
\item[Easy 2.2] 61.2\,mW ok; 250\,mW at the line; 400\,mW over.
\item[Easy 2.3] Design current to the typ/$V_F$ band; the 30\,mA abs-max is the cliff edge you never plan to visit.
\item[Easy 2.4] R1 between nodes 3--7 (470\,$\Omega$); D2 between 7--0; node 0 is ground.
\item[Med (derive)] $P=VI$, $V=IR \Rightarrow P=(IR)I=I^2R$; $I=V/R \Rightarrow P=V\cdot V/R=V^2/R$; solving $V^2/R \le P$ for $R$ gives $R \ge V^2/P$ — careful: the \emph{minimum} resistance is set by the rating, not the maximum.
\item[Med (220)] $I_{\mathrm{nom}} = 3/220 = 13.6$\,mA; band $\approx 11.0$--$16.6$\,mA; still good, slightly dimmer.
\item[Med (1k+1k)] $v(2)=4.5$\,V, $i=4.5$\,mA — a divider preview.
\item[Med (colours)] 220\,$\Omega$: 209--231\,$\Omega$; yellow-violet-red-gold $=4\,7\times10^2=4.7$\,k$\Omega$: 4465--4935\,$\Omega$.
\item[Hard (two LEDs)] With IS mismatched by 10\%, the hog takes $\sim$2$\times$ the other's current — and the gap widens with heat. Per-LED resistors is the fix.
\item[Hard (design)] $I_{\max}$ occurs at 5\,V, $V_F$=1.8, $R_{\min}$: need $R\ge(5.25-1.8)/0.02=172.5\,\Omega$; $I_{\min}$ at 3.135\,V, $V_F$=2.2, $R_{\max}$: need $R\le0.935/0.01=93.5\,\Omega$. Impossible as stated — the answer is a smaller tolerance or two supply rails; the exercise is the proof.
\end{description}

\chapternotes

\begin{enumerate}
\item T.~R. Kuphaldt, \emph{Lessons in Electric Circuits}, Vol.~I ``DC'',
chapters 2 (Ohm's law) and 1.6 (power): the $P=VI$ forms and the ratings
habit; its own SPICE introduction validates our netlist-first order
\cite{aacvol1}.
\item OpenStax, \emph{University Physics} Vol.~2: the derivation of
$V=IR$ for ohmic materials and where non-ohmic devices break it —
Section~\ref{sec:nonlinear}'s correct statement \cite{openstaxv2}.
\item ngspice manual, \texttt{.op}/\texttt{.dc} analysis and \texttt{.control}
scripting: the exact directives of Listings~\ref{lst:sweepcir}; version 47
pinned in \texttt{code/README.md} \cite{ngspice}.
\item The E24 preferred-value series and colour code are industry standard;
the decoding drill mirrors Table~1.1 \cite{aacvol1}.
\item The LED model parameters (\texttt{IS}, \texttt{N}) are chosen so the
simulated knee lands near a real red LED's $V_F$; they are teaching values,
not a vendor's characterised model — the honesty point of
Example~\ref{ex:threeways}.
\end{enumerate}

\begin{rulebox}[Watch instead]
\begin{itemize}
\item \emph{Introduction to circuits and Ohm's law} (Khan Academy, 11:39) --- $V = IR$ and power in a simple loop. \url{https://youtu.be/NYQH1nK9PQ4} Listed in Book 8 (ch12), accessed 2026-10-11.
\end{itemize}
\end{rulebox}

